\documentclass[11pt]{article}

\usepackage[margin=1in]{geometry}
\usepackage{amsmath,amssymb,booktabs,array,enumitem,hyperref,xcolor}
\usepackage[T1]{fontenc}
\usepackage{lmodern}
\hypersetup{colorlinks=true,linkcolor=blue,urlcolor=blue,citecolor=blue}

\newcommand{\ppocf}{\textsc{PPO-CF}}
\newcommand{\pposca}{\textsc{PPO-SCA}}
\newcommand{\ppo}{\textsc{PPO}}
\newcommand{\todo}[1]{\textcolor{red}{\textbf{TODO:} #1}}

\title{Steering Policy Optimization with Sparse Counterfactual Advantages}
\author{Paper pitch and writing blueprint}
\date{}

\begin{document}
\maketitle

\begin{abstract}
Policy-gradient agents are usually controlled through rewards, although the
policy is updated through an advantage estimator. We study \emph{counterfactual
advantage steering}: augmenting PPO's actor update with a policy-centered,
action-selective counterfactual advantage while leaving the environment reward,
GAE targets, and critic update unchanged. We first show that a state-only
advantage perturbation cancels from the expected policy gradient; behavioral
control therefore requires an action-selective residual. Guided by this
observation, our method queries exact local counterfactual action values at a
small subset of rollout states and injects the resulting centered landscape
into PPO's sampled-action objective. Across long-horizon visual control tasks,
sparse counterfactual queries consistently improve learning relative to PPO and
often match or exceed a more expensive direct counterfactual baseline. In
contrast, globally propagating sparse labels through an imperfect distilled
landscape can be unreliable. These results identify sparse, action-aligned
counterfactual advantages as a lightweight control interface for policy
optimization.
\end{abstract}

\section*{Abstract 2: problem--need--method--importance}
\begin{abstract}
PPO is typically guided by modifying rewards, but reward design is often
expensive, changes the task being optimized, and can unintentionally alter the
critic's learning target. This is especially limiting in structured simulators,
where a small number of local counterfactual evaluations may reveal which action
is preferable without requiring a new reward function. We ask whether this
privileged but sparse information can be used to guide the policy update
directly. We introduce \emph{counterfactual advantage steering}, which adds a
policy-centered, action-selective counterfactual advantage to PPO's actor loss
at a small subset of rollout states, while preserving the environment reward,
GAE returns, and critic update. Our analysis shows why action selectivity is
necessary: state-only advantage offsets cancel from the expected policy
gradient. Empirically, sparse exact counterfactual queries improve learning on
long-horizon visual control tasks and are often more reliable than applying a
distilled counterfactual estimate everywhere. By separating policy guidance
from reward and value-target design, this work provides a practical route for
using structured simulators to make policy optimization more efficient and
controllable.
\end{abstract}

\section*{One-sentence paper claim}
\begin{quote}
\emph{A small number of exact, policy-centered counterfactual advantages can
steer PPO's actor update more effectively than either ordinary PPO or a noisy
dense approximation, without modifying rewards or critic targets.}
\end{quote}

\paragraph{Scope.}
This is deliberately a paper about \emph{sparse exact counterfactual advantage
queries}, not about TD clipping, generic advantage bias, a learned dense
landscape, or reward shaping. The secondary-goal work supports the intuition
that the actor advantage is a control interface, but it is not a primary
empirical claim in this version of the paper.

\section{Introduction: why this problem matters}
\label{sec:intro}

\paragraph{Opening.}
Start with the mismatch between how PPO is commonly controlled and how it is
actually optimized. Rewards and critic targets are slow, global control
mechanisms. The actor, however, only sees a scalar advantage associated with
the sampled action. A simulator or structured environment can sometimes answer
the local question ``what would happen under the other actions at this state?''
The question is whether a small number of such local answers can improve policy
optimization without rewriting the reward function.

\paragraph{Gap.}
Full counterfactual PPO is computationally expensive because it needs an
all-action landscape at many states. An obvious alternative is to learn a
landscape from sparse labels and apply it everywhere. The experiments show that
this alternative is not automatically safe: approximation error in the
action-relevant direction can hurt the policy even when the student is trained
on the same label budget.

\paragraph{Contribution paragraph to write.}
State exactly three contributions.
\begin{enumerate}[leftmargin=1.6em]
  \item We isolate the action-selective component of an advantage perturbation
  as the part that affects the expected actor gradient.
  \item We introduce \pposca{}: PPO with sparse exact counterfactual advantage
  queries, injected into the ordinary sampled-action actor loss while rewards,
  returns, and critic targets are untouched.
  \item We show on four main MiniGrid tasks that sparse exact queries improve
  learning over PPO, compare favorably with direct \ppocf{}, and are more
  reliable than naively applying an imperfect distilled landscape everywhere.
\end{enumerate}

\paragraph{Do not write in the introduction.}
Do not lead with TD clipping, adversarial attacks, backdoors, uncertainty
selection, active distillation, or the custom secondary-goal environment.
Those are either separate projects or ablations which obscure the central
algorithmic claim.

\section{Background and mechanism: advantage perturbations}
\label{sec:mechanism}

\paragraph{Core derivation.}
Let an actor use an altered estimator
\begin{equation}
  \widehat A^f(s,a) = \widehat A(s,a) + f(s,a).
\end{equation}
The additional expected policy-gradient contribution is determined by
\begin{equation}
  \Delta_f(s,a)
  = f(s,a) - \mathbb{E}_{a'\sim\pi(\cdot\mid s)}[f(s,a')].
  \label{eq:centered-residual}
\end{equation}
This equation should be the central conceptual figure or boxed proposition:
the state-only portion of $f$ is removed by the score-function identity, while
the policy-centered action-dependent residual remains.

\paragraph{Result to use from the advantage-perturbation folder.}
Use the \emph{mechanistic control}, not TD clipping. The relevant experiment is
the matched state-only versus action-selective perturbation result from
\texttt{Adv\_perturbation\_experiments/TD\_CLIPPING.zip}. It establishes the
following qualitative fact: at the same realized perturbation scale, a
state-varying but action-independent term is an exact null for the expected
actor gradient, whereas an action-selective term changes learning. Present it
as a sanity check for Equation~\eqref{eq:centered-residual}, not as a second
main benchmark suite.

\paragraph{Suggested figure.}
\begin{center}
\fbox{\parbox{0.85\linewidth}{
\textbf{Figure 1: Action selectivity is necessary.} Plot performance or update
displacement against perturbation strength for clean, state-only ($\beta=0$),
and action-selective ($\beta=1$) perturbations. Hold realized perturbation
scale fixed. The visual message is ``same magnitude, different effect.''}}
\end{center}

\section{Method: sparse counterfactual advantage queries}
\label{sec:method}

\paragraph{Counterfactual landscape.}
At a queried simulator state $s$, evaluate local action values
$Q_{\rm CF}(s,a)$ and policy-center them:
\begin{equation}
 A_{\rm CF}(s,a) = Q_{\rm CF}(s,a) -
 \sum_{a'}\pi_{\rm old}(a'\mid s)Q_{\rm CF}(s,a').
 \label{eq:cf-advantage}
\end{equation}
The centering is essential: it makes the intervention explicitly
action-selective and prevents a state-wide offset from being mistaken for
useful policy guidance.

\paragraph{PPO-SCA update.}
For a queried transition, modify only the sampled-action actor advantage:
\begin{equation}
 A_{\rm actor,t} = \operatorname{norm}(A_{\rm GAE,t}) +
 \beta\,\frac{A_{\rm CF}(s_t,a_t)}{\sigma_{\rm GAE}}.
 \label{eq:sca-update}
\end{equation}
For unqueried transitions, use standard normalized GAE. In prose, explicitly
state the three invariants:
\begin{enumerate}[leftmargin=1.6em]
  \item The environment reward is unchanged.
  \item GAE returns and value targets are unchanged.
  \item The counterfactual term is detached and affects only the actor loss.
\end{enumerate}
This distinguishes the method from reward shaping, critic-target shaping, and
an auxiliary value function.

\paragraph{Query budget.}
The main setting uses exact all-action counterfactual landscapes at 2\% of
rollout states and uses $\beta=0.75$. Before submission, write down a fixed
selection rule for $\beta$ (or retain the full $\{0.25,0.75,1.5\}$ sweep in the
appendix). The paper must not imply that $\beta=0.75$ was uniquely selected
after inspecting only the best curves.

\paragraph{Pseudocode.}
Include a short algorithm box: collect PPO rollout; uniformly choose 2\% of
states; restore each selected simulator state and evaluate actions; form
Equation~\eqref{eq:cf-advantage}; add the detached sampled-action term in
Equation~\eqref{eq:sca-update}; run ordinary PPO update. Keep this under half a
page.

\section{Experimental protocol}
\label{sec:protocol}

\subsection{Algorithms retained in the main paper}
Keep only the following three learning curves in the main figures.
\begin{center}
\begin{tabular}{p{0.24\linewidth}p{0.66\linewidth}}
\toprule
\textbf{Method} & \textbf{Purpose} \\
\midrule
\ppo{} (GAE) & Standard policy-optimization baseline. \\
Direct \ppocf{} & Counterfactual baseline using its direct all-action objective;
shows what sparse queries are competing with. \\
\pposca{} (2\% exact queries) & Proposed method: exact centered landscape added to
the sampled-action actor advantage. \\
\bottomrule
\end{tabular}
\end{center}

\paragraph{Methods removed from the main comparison.}
Remove the following from the title, abstract, main figures, and headline table:
uniform distillation, uncertainty distillation, active distillation, multiple
student-selection schemes, TD clipping, shuffled/attack variants, and manual
secondary-goal steering. They create an unreadable six-method plot and weaken
the paper's causal message.

\paragraph{Where the removed methods go.}
Put uniform distillation in an ablation subsection because it motivates using
exact sparse labels directly. Put active/uncertainty selection in the appendix
as preliminary evidence only; it currently has three seeds. Put the
secondary-goal result in the appendix or a final qualitative discussion. Do not
silently delete negative results: show all removed variants in an appendix
table with their configurations and seeds.

\subsection{Environments and selection}
The strongest four-panel learning-curve figure from
\texttt{notebooks/06\_plot\_runs\_eval\_training\_curves.ipynb} is:
\begin{enumerate}[leftmargin=1.6em]
  \item MiniGrid-RedBlueDoors-6x6-v0;
  \item MiniGrid-KeyCorridorS3R3-v0;
  \item MiniGrid-DoorKey-6x6-v0;
  \item MiniGrid-UnlockPickup-v0.
\end{enumerate}
KeyCorridor S3R2 should be in the appendix because its outcomes are visibly
more variable. The pre-specified E2/E3 panel also includes Taxi. Re-export Taxi
in the same evaluation format before submission so that the formal four-task
table can use Taxi, DoorKey, UnlockPickup, and RedBlueDoors without mixing
evaluation pipelines.

\paragraph{Selection disclosure.}
The current four visual environments were selected after inspecting existing
curves. Do not make a post-selected average the primary inferential result.
Either (i) freeze the proposed four-task set and run fresh confirmatory seeds,
or (ii) report every completed task, including KeyCorridor S3R2, in the main
aggregate and reserve the four-panel figure for readability. This is important
for an ICLR submission.

\subsection{Evaluation and reporting}
Report mean success rate and episodic return as training curves with five-seed
mean and uncertainty bands. Use time-to-threshold (e.g., frames to 90\% success)
and final evaluation success in the main table. A normalized area-under-curve
statistic can be reported in the table or appendix as a summary, but it should
not be the centerpiece of the abstract.

\section{Results}
\label{sec:results}

\subsection{Main result: sparse exact counterfactual queries improve learning}

\paragraph{What to write.}
Open this section with the qualitative conclusion, then point to the curves:
``Across the displayed tasks, sparse exact counterfactual queries improve or
accelerate PPO learning. The largest gains occur in tasks where the policy must
make action-critical sequential decisions, rather than merely discover a
terminal reward.'' Avoid saying the method universally dominates direct
\ppocf{}; DoorKey is an explicit counterexample.

\begin{center}
\begin{tabular}{lcccc}
\toprule
& \multicolumn{3}{c}{Normalized success during training} &
Frames to 90\% success \\
\cmidrule(lr){2-4}
\textbf{Environment} & \ppo{} & Direct \ppocf{} & \pposca{} &
\ppo{} $\rightarrow$ \pposca{} \\
\midrule
RedBlueDoors 6x6 & 0.528 & 0.636 & \textbf{0.710} & 848k $\rightarrow$ 695k \\
KeyCorridor S3R3 & 0.677 & 0.655 & \textbf{0.826} & 268k $\rightarrow$ 152k \\
DoorKey 6x6 & 0.492 & \textbf{0.633} & 0.534 & 285k $\rightarrow$ 273k \\
UnlockPickup & 0.588 & 0.656 & \textbf{0.754} & 653k $\rightarrow$ 404k \\
\bottomrule
\end{tabular}
\end{center}

\paragraph{Table note to include in the final paper.}
The normalized-success column above is a working extraction from the five-seed
curves in Notebook 06, not yet a locked camera-ready statistic. Recompute it in
one evaluation script, record uncertainty intervals, and state the exact
integration convention. The numbers are useful for selecting the story but
must be reproducible before being treated as final paper values.

\paragraph{Figure 2.}
Use a four-panel success-rate learning curve: PPO in blue, direct \ppocf{} in
orange, and \pposca{} in green. Do not include active distillation. The existing
focused plot can be simplified by removing its black active-distilled line.

\subsection{Sparse exact labels are safer than naive dense distillation}

\paragraph{What to write.}
This is the key ablation, not a co-equal method. The question is whether the
same 2\% label budget should be used immediately at queried states or learned
and broadcast to all states. Write: ``The benefit is not explained simply by
adding a dense auxiliary signal. On several tasks, approximation error in a
student landscape is amplified when it is applied at every state, whereas exact
queried values retain the correct local action ranking.''

\paragraph{Results to use.}
At $\beta=0.75$, exact queries have higher recorded success during training than
uniform distillation in RedBlueDoors, Taxi, and UnlockPickup. The clearest
example is RedBlueDoors: the uniform student variant has much lower final
success and learning progress than the exact queried variant. DoorKey is the
important exception, where distillation is competitive or better; report this
plainly.

\paragraph{Suggested figure/table.}
Use one compact table with Query, Uniform-distill, and the difference for the
pre-specified E2 environments. Keep the complete beta sweep in the appendix.
Do not add active or uncertainty variants to this figure.

\subsection{Action selectivity is the operative ingredient}

\paragraph{What to write.}
Tie the perturbation experiment back to the method: ``The centering in
Equation~\eqref{eq:cf-advantage} is not cosmetic. Matched perturbations show
that state-dependent offsets do not produce the observed behavior change;
action alignment is required.'' Include the $\beta=0$ identity/null and, if
available, an action-permuted counterfactual landscape control. This is the
section that makes the paper more than an empirical oracle heuristic.

\subsection{Cost and scope}

\paragraph{What to write.}
Report counterfactual branches, wall-clock time, and query fraction for PPO,
direct \ppocf{}, and \pposca{}. Do not imply that the oracle is free. State that
the present study assumes a simulator state-restoration interface and
deterministic/local counterfactual rollouts. This makes the contribution an
algorithmic result for structured simulators, not a claim about arbitrary
model-free environments.

\section{Discussion and limitations}
\label{sec:discussion}

\paragraph{Main discussion.}
The results suggest that the useful resource is not a generic dense learned
bonus. It is a correctly centered, action-level signal at states where an exact
counterfactual can be trusted. Explain that this produces a practical spectrum:
ordinary PPO uses no simulator intervention; direct \ppocf{} uses more complete
counterfactual supervision; \pposca{} obtains much of the benefit from sparse
queries.

\paragraph{Limitations to state explicitly.}
\begin{enumerate}[leftmargin=1.6em]
  \item The current evidence is MiniGrid-heavy; a final submission should add
  at least one non-MiniGrid environment with a valid restoration oracle.
  \item The oracle is privileged simulator access. This is not a model-free
  method unless a learned world model replaces the oracle.
  \item The counterfactual horizon and $\beta$ are task-sensitive.
  \item Distillation is not yet a solved scaling mechanism; the reported
  negative result is part of the conclusion.
\end{enumerate}

\paragraph{Secondary-goal work.}
Mention only as a future direction or appendix demonstration: an actor-only,
reward-free advantage term can induce blue-token collection in a custom
DetourPickup environment while retaining goal completion in the available
two-seed experiment. Do \emph{not} claim a robust primary--secondary trade-off
or general controllability from this result.

\section{Appendix plan}
\begin{enumerate}[leftmargin=1.6em]
  \item Full five-environment Notebook 06 curves, including KeyCorridor S3R2.
  \item Full E2 beta sweep for exact queried and uniform distilled landscapes.
  \item E3 active/uncertainty results, marked as three-seed exploratory results.
  \item Detailed derivation of the centered residual and the state-only null.
  \item Oracle restoration checks, centering diagnostics, query counts, and
  computational cost.
  \item All configurations, seeds, evaluation seeds, and selection rules.
  \item Secondary-goal implementation and its two-seed qualitative result.
\end{enumerate}

\section*{Final checklist before converting this pitch into a submission}
\begin{enumerate}[leftmargin=1.6em]
  \item Recreate every headline number with one locked evaluation script.
  \item Freeze the environment and $\beta$ selection rules; run confirmation
  seeds if the four-task figure remains post-selected.
  \item Add Taxi to the common evaluation/plot pipeline or omit it consistently.
  \item Add a matched action-permutation control if it is not already complete.
  \item Measure query/oracle wall-clock cost against direct \ppocf{}.
  \item Add one non-MiniGrid task before treating the work as ICLR-ready.
\end{enumerate}

\end{document}
